\textbf{Abstrak}

Robot bergerak otonom yang beroperasi pada lingkungan indoor umumnya mengandalkan sensor LiDAR dua dimensi untuk melakukan deteksi obstacle dan navigasi. Meskipun memiliki akurasi pengukuran jarak yang tinggi, LiDAR 2D hanya melakukan pemindaian pada satu bidang horizontal sehingga obstacle yang berada di luar bidang tersebut berpotensi tidak terdeteksi. Kondisi ini menimbulkan area \textit{blind-spot} yang dapat meningkatkan risiko tabrakan, khususnya terhadap obstacle berprofil rendah yang berada di dekat permukaan lantai. Selain itu, penggunaan parameter navigasi yang bersifat tetap pada metode \textit{Artificial Potential Field} (APF) sering kali menghasilkan perilaku navigasi yang kurang adaptif terhadap perubahan kondisi lingkungan.

Penelitian ini mengusulkan sistem navigasi robot beroda otonom berbasis logika fuzzy dengan fusi sensor LiDAR dan kamera termal untuk meningkatkan keselamatan navigasi pada lingkungan indoor. Sistem yang dikembangkan mengintegrasikan data LiDAR 2D dan citra termal melalui pendekatan estimasi traversability. Citra termal diproses untuk menghasilkan informasi traversability termal dan \textit{virtual LaserScan} yang merepresentasikan obstacle pada area yang tidak teramati oleh LiDAR. Selanjutnya, traversability LiDAR dan traversability termal digabungkan menjadi traversability efektif yang digunakan sebagai masukan sistem navigasi. Untuk meningkatkan kemampuan adaptasi terhadap kondisi lingkungan, penelitian ini menerapkan pengendali fuzzy Mamdani yang mengatur parameter \textit{speed scale} dan \textit{repulsive gain} pada APF berdasarkan jarak obstacle terdekat dan nilai traversability efektif.

Pengujian dilakukan pada lingkungan simulasi ROS Noetic dan Gazebo menggunakan delapan skenario yang terdiri atas empat konfigurasi baseline dan empat konfigurasi fuzzy. Setiap skenario dijalankan sebanyak sepuluh kali percobaan independen sehingga diperoleh total delapan puluh \textit{trial}. Evaluasi dilakukan menggunakan metrik \textit{Success Rate}, \textit{Collision-Free Rate}, \textit{minimum clearance}, waktu tempuh, panjang lintasan, serta karakteristik pergerakan robot. Analisis statistik dilakukan menggunakan uji Mann--Whitney U dan \textit{rank-biserial correlation}.

Hasil pengujian menunjukkan bahwa integrasi kamera termal dan pengendali fuzzy meningkatkan performa navigasi terutama pada lingkungan yang mengandung obstacle berprofil rendah dan obstacle dinamis. Konfigurasi fuzzy mampu mempertahankan \textit{Success Rate} dan \textit{Collision-Free Rate} sebesar 100% pada seluruh skenario pengujian, serta menghasilkan peningkatan \textit{minimum clearance} yang signifikan dibandingkan konfigurasi baseline. Peningkatan keselamatan tersebut diperoleh dengan konsekuensi berupa peningkatan waktu tempuh yang relatif moderat akibat perilaku navigasi yang lebih konservatif. Selain itu, penggunaan kamera termal terbukti mampu mengurangi dampak area \textit{blind-spot} LiDAR melalui deteksi obstacle yang tidak sepenuhnya teramati oleh sensor laser.

Hasil penelitian menunjukkan bahwa fusi sensor LiDAR dan kamera termal yang dipadukan dengan pengendali fuzzy Mamdani mampu meningkatkan keselamatan navigasi robot beroda otonom tanpa mengubah struktur dasar ROS Navigation Stack. Pendekatan yang diusulkan berpotensi menjadi solusi untuk meningkatkan robustitas sistem navigasi robot pada lingkungan indoor yang mengandung obstacle berprofil rendah dan kondisi operasi yang dinamis.

\textbf{Kata kunci}: robot bergerak otonom, LiDAR 2D, kamera termal, fusi sensor, traversability, logika fuzzy Mamdani, Artificial Potential Field, navigasi robot.
